% 01-introduction.tex
\section{Introduction}
\label{sec:intro}

Service meshes have become the default way to add mutual TLS, fine-grained
traffic policy, and request-level telemetry to microservices on Kubernetes,
without changing application code. The price is paid on every request: traffic
that used to flow directly between two pods now crosses one or more proxies,
each of which terminates connections, encrypts and decrypts traffic, and often
parses application-layer protocols. In deep call graphs, this cost is paid at
every hop, and it shows up where operators care most: in the tail of the
latency distribution and in the CPU that must be provisioned for it.

The design space of service meshes has changed considerably in the last few
years. The original \emph{sidecar} architecture, in which a proxy runs next to
every application container (Istio with Envoy, Linkerd with its Rust
\texttt{linkerd2-proxy}), now competes with \emph{sidecar-less} designs. Istio
Ambient splits the data plane into a per-node L4 proxy (ztunnel) and optional,
shared L7 proxies (waypoints), while Cilium implements L3/L4 functions in eBPF
inside the kernel and delegates L7 processing to a per-node Envoy. Each design
moves work to a different place (the pod, the node, or the kernel) and changes
how many proxies a request traverses, so their costs cannot be inferred from
one another.

The most detailed study of mesh overheads to date, MeshInsight~\cite{zhu2023meshinsight},
decomposed the cost of the Istio sidecar into its components and found that it
can raise latency by up to 269\% and CPU usage by up to 163\%. That analysis,
however, covers only the Envoy sidecar, profiles it with at most one
outstanding request, and its model explicitly cannot predict tail latency. More recent comparisons of sidecar and
sidecar-less meshes~\cite{bremlerbarr2024mtls} examine mTLS alone, on a
synthetic server with an artificial delay, without L7 processing, waypoints,
energy, or a predictive model. Operators choosing a mesh today therefore lack a
same-testbed comparison of the architectures they can actually deploy.

This paper provides that comparison. We measure seven configurations on the same
testbed: no mesh, Istio sidecar, Istio Ambient with and without a waypoint,
Linkerd, and Cilium service mesh with and without L7 policy, each with east-west encryption enabled
in its default data path: workload-level mutual TLS for Istio and Linkerd,
and node-to-node WireGuard for Cilium. We drive an echo
microbenchmark and two widely used microservice benchmarks, Online
Boutique~\cite{onlineboutique} and Hotel Reservation from
DeathStarBench~\cite{gan2019deathstarbench}, with open-loop load, which avoids
the coordinated-omission bias that hides tail latency~\cite{tene2015latency}.
Our contributions are:

\begin{itemize}
  \item \textbf{A same-testbed measurement of seven mesh configurations}
        (RQ1). We report p99 and p99.9 latency, maximum SLO-compliant
        throughput, CPU, memory, and energy per request, with ten randomized
        repetitions per cell for latency and resources (five for Cilium
        without L7 policy), three for energy, and
        non-parametric statistics. Sidecar-based
        Istio sustains only 30\% of the no-mesh throughput on the echo
        benchmark, while ambient L4 mode retains 78\%.
  \item \textbf{An attribution of overhead beyond Envoy} (RQ2). Using perf and
        eBPF, we break the cost down into TLS, L7 parsing, socket I/O, context
        switches, and kernel time, including ztunnel, the Ambient waypoint,
        \texttt{linkerd2-proxy}, and Cilium's agent, per-node Envoy, and
        WireGuard. Ablations show that Cilium's cost comes from the L7
        redirect and Hubble flow logging, not from its eBPF programs.
  \item \textbf{A model of the tail and its limits} (RQ3). A compositional
        model of per-hop proxy cost, cross-node hops, and queueing, fitted on
        two applications, explains the p99 of the third within 11--13\% for
        the L4 data planes; for per-pod sidecars, queueing at the busiest
        proxy instance cuts the error from 82\% to 19\% (Linkerd) and from
        135\% to 52\% (Istio), while node utilization does not explain it.
        Node utilization bounds throughput within 19\% on the two
        microservice applications. A placement experiment shows that node
        load explains the tail of L7 data planes only partly.
  \item \textbf{A decision guide} (RQ4) that maps latency, CPU, and energy
        budgets to the mesh that satisfies them.
\end{itemize}

We also partially reproduce MeshInsight: the CPU overhead of the Istio
sidecar on Online Boutique falls within the published range, but its latency
overhead is larger (\cref{sec:validation}). We release the testbed, scripts, and data to
support reproduction at \anonrepo.

The rest of the paper is organized as follows. \Cref{sec:background} describes
the mesh architectures we compare, \cref{sec:related} discusses related work,
\cref{sec:method} presents the testbed and methodology, and
\cref{sec:rq1,sec:rq2,sec:rq3,sec:rq4} answer the four research questions.
\Cref{sec:threats} discusses threats to validity and \cref{sec:conclusion}
concludes.
